\section{补充：二阶优化方法}

\begin{figure}[H]
\centering
\includegraphics[width=0.95\textwidth]{slides-images/slide-058.jpg}
\caption{二阶优化：利用 Hessian 矩阵直接跳到极值点\protect\footnotemark}
\end{figure}
\footnotetext{Slides 第62页。}

前面介绍的所有方法都是\textbf{一阶优化}——只用到梯度（一阶导数）信息。\textbf{二阶优化}方法（如 Newton 方法）利用 Hessian 矩阵（二阶导数矩阵）直接估计极值点位置：

$$W^* = W - H^{-1} \nabla_W L$$

\begin{warningbox}{二阶方法在深度学习中不实用}
对于有 $N$ 个参数的模型，Hessian 矩阵大小为 $N \times N$。现代大模型有数十亿参数，存储和求逆 Hessian 是完全不可行的。因此深度学习几乎全部使用一阶方法（SGD、Adam 等）。二阶方法在小规模凸优化问题中仍然有用。
\end{warningbox}

%% ========== 总结与延伸 ==========

